% ===== Preâmbulo compartilhado (Tectonic / XeTeX) — documentos da FTL094 =====
\documentclass[11pt,a4paper]{article}
\usepackage{fontspec}
\usepackage[brazil]{babel}
\usepackage{geometry}
\geometry{a4paper,margin=2.5cm}
\usepackage{amsmath,amssymb}
\usepackage{graphicx}
\usepackage{booktabs}
\usepackage{array}
\usepackage{longtable}
\usepackage{tabularx}
\usepackage{float}
\usepackage{enumitem}
\usepackage{xcolor}
\usepackage{tikz}
\usetikzlibrary{arrows.meta,positioning,automata,shapes.geometric,fit,backgrounds,calc}
\usepackage{pgfplots}
\pgfplotsset{compat=1.18}
\usepackage{listings}
\usepackage{caption}
\usepackage{fancyhdr}
\usepackage[hidelinks]{hyperref}

% ----- cores -----
\definecolor{ufamblue}{RGB}{0,51,102}
\definecolor{codegreen}{rgb}{0,0.5,0}
\definecolor{codeblue}{rgb}{0,0,0.7}
\definecolor{codered}{rgb}{0.6,0,0}
\definecolor{lgray}{rgb}{0.95,0.95,0.95}

% ----- listagens C -----
\lstdefinestyle{c}{language=C,basicstyle=\ttfamily\footnotesize,
  keywordstyle=\color{codeblue}\bfseries,commentstyle=\color{codegreen}\itshape,
  stringstyle=\color{codered},numbers=left,numberstyle=\tiny\color{gray},
  frame=single,breaklines=true,tabsize=2,showstringspaces=false,columns=fullflexible}

% ----- constantes do projeto -----
\newcommand{\projnome}{Sistema de Navegação Autônoma para o Robô Móvel Khepera~IV}
\newcommand{\disciplina}{FTL094 --- Engenharia de Software de Tempo Real}
\newcommand{\periodo}{Período Letivo 2026/1}
\newcommand{\versaodoc}{1.0}

% ----- cabeçalho/rodapé -----
\pagestyle{fancy}\fancyhf{}
\rhead{\small \tipodoc}
\lhead{\small Khepera~IV --- FTL094}
\cfoot{\thepage}
\renewcommand{\headrulewidth}{0.4pt}

% ----- coluna X centralizada p/ tabularx -----
\newcolumntype{C}{>{\centering\arraybackslash}X}

% ----- capa padronizada -----
\newcommand{\fazcapa}{%
\begin{titlepage}\centering
{\large\bfseries UNIVERSIDADE FEDERAL DO AMAZONAS (UFAM)}\\[3pt]
{\normalsize Centro de P\&D em Tecnologia Eletrônica e da Informação --- CETELI}\\[2pt]
{\normalsize Programa de Pós-Graduação em Engenharia Elétrica --- PPGEE}\\[26pt]
{\large \disciplina}\\[2pt]
{\normalsize \periodo}\\[64pt]
\rule{\textwidth}{1.2pt}\\[12pt]
{\Huge\bfseries\color{ufamblue} \tipodoc}\\[14pt]
{\Large \projnome}\\[10pt]
\rule{\textwidth}{1.2pt}\\[48pt]
{\large\bfseries Equipe}\\[8pt]
{\normalsize Gabriel Henrique de Souza Nazaré}\\[3pt]
{\normalsize Luan Alexandre Ramos Barreto}\\[3pt]
{\normalsize Matheus Rocha Canto}\\[3pt]
{\normalsize Matheus Serrão Uchôa}\\[3pt]
{\normalsize Ricardo Mendonça Braz}\\[3pt]
{\normalsize Yan Fernandes da Silva}\\[40pt]
{\normalsize Docente responsável: Prof. Dr.~[nome do docente]}\\
\vfill
{\small Documento de tipo \emph{\tipodoc} --- Versão \versaodoc}\\[2pt]
{\small Manaus --- AM, \datadoc}
\end{titlepage}}

% ----- tabela de controle de versões -----
\newcommand{\controleversoes}{%
\section*{Controle de versões}
\begin{tabularx}{\textwidth}{@{}l l l X@{}}
\toprule
\textbf{Versão} & \textbf{Data} & \textbf{Autor} & \textbf{Descrição} \\
\midrule
1.0 & \datadoc & Equipe FTL094 & Emissão inicial do documento. \\
\bottomrule
\end{tabularx}
\bigskip}
